% =============================================================================
% Ready-to-paste LaTeX: the evaluation protocol and the metrics.
% Answers supervisor request #2 -- "explain the process of the evaluation, and which
% metrics produced the feature-extraction accuracy".
%
% Placement: the Methods chapter, as its own section immediately before the Results.
% Every number quoted here is reproduced by experiments/evaluation/exp_statistical_power.py
% (results/statistical_power.json), experiments/emotion/exp_emotion_regression.py and
% experiments/cross_dataset/exp_zero_shot.py (results/zero_shot.json).
%
% CITATION NOTE: uses nadeau2003inference (BibTeX entry given in statistical_power.tex)
% and the standard multi-label metric references. If the following are not yet in
% bib/library.bib, add:
%
%   @article{sokolova2009systematic,
%     title   = {A Systematic Analysis of Performance Measures for Classification Tasks},
%     author  = {Sokolova, Marina and Lapalme, Guy},
%     journal = {Information Processing \& Management},
%     volume  = {45}, number = {4}, pages = {427--437}, year = {2009},
%     doi     = {10.1016/j.ipm.2009.03.002}
%   }
% =============================================================================

\section{Evaluation Protocol}
\label{sec:evaluation_protocol}

Because this thesis compares a two-stage pipeline against a direct one, the evaluation
has to answer two different questions with two different kinds of metric, and it has to
do so under a resampling scheme that neither flatters nor penalises either architecture.
This section states the protocol in full: what is measured at each stage, which metric is
used and why, how the data are split, how hyperparameters are chosen, and how differences
between systems are tested for significance.

\subsection{What is evaluated, and with which metric}

The pipeline decomposes into two stages that are evaluated separately and then
end-to-end.

\paragraph{Stage 1: audio $\rightarrow$ emotion (regression).}
Each audio representation is judged by how well it predicts the eight human emotion
ratings. Since the ratings are continuous values on a 1--9 scale, this is a regression
problem, and the metric is the coefficient of determination $R^2$, computed
\emph{separately for each of the eight emotions} and then averaged. $R^2$ is reported
rather than a correlation because it is defined against the variance of the target and is
therefore interpretable as the proportion of rating variance the audio features recover:
$R^2 = 0$ is the performance of always predicting the training mean, and negative values
are possible and meaningful. The root mean squared error (RMSE) is reported alongside it,
expressed in rating points, because $R^2$ alone does not convey whether an error is
practically large; an RMSE of $1.0$ means the model is typically one scale point away.

\textbf{This mean $R^2$ is the metric referred to as ``feature-extraction accuracy''.} It
is what ranks the five audio representations against one another
(Table~\ref{tab:feature_ranking} in Chapter~\ref{ch:evaluation}), because it measures
exactly the property the pipeline needs from a representation: how much affective
information it carries. The genre metrics below cannot serve this purpose, since they
also confound the quality of the second-stage classifier.

\paragraph{Stage 2: emotion $\rightarrow$ genre (multi-label classification).}
A film belongs to several genres at once, so genre prediction is multi-label: the target
is a binary vector, and the model may be partly right. Three metrics are computed, and
they are not interchangeable:

\begin{itemize}
  \item \textbf{Macro-F1} -- the F1 score computed per genre and then averaged with equal
    weight per genre. This is the headline metric.
  \item \textbf{Exact-match ratio} (subset accuracy) -- the fraction of clips whose entire
    genre vector is predicted correctly. Strict, and reported for completeness.
  \item \textbf{Hamming loss} -- the fraction of individual label decisions that are wrong.
\end{itemize}

Macro-F1 is the headline because the genre distribution is severely imbalanced (Drama
appears on 239 of 346 clips, Documentary on 17). Under such imbalance the other two
metrics reward a model for ignoring the rare genres: a trivial classifier that always
predicts the most frequent label set attains a \emph{better} exact-match ratio and a
\emph{better} Hamming loss than every real model tested here, while being useless. Macro
averaging is the standard remedy \citep{sokolova2009systematic}, since it gives a genre
with 17 positives the same weight as one with 239 and therefore cannot be gamed by
neglecting the minority classes. All three numbers are reported in the results tables so
that this behaviour is visible rather than hidden.

For the cross-corpus comparison with \citet{ma2021computational}, macro-averaged
precision and recall are reported next to macro-F1, matching their reporting format: F1
alone cannot distinguish a model that over-predicts (recall high, precision low) from one
that under-predicts, and the two failure modes call for different remedies.

\paragraph{Sanity check.}
Before either stage is interpreted, the extraction pipeline is verified end-to-end on a
deliberately easy auxiliary task: predicting the corpus's own 12 balanced discrete emotion
categories from the embedding. Because the classes are balanced and single-label, plain
accuracy is appropriate there, and chance is exactly $1/12 = 0.083$. Scoring well above
chance confirms that the cached embeddings are aligned with their labels and carry
signal --- a check that would have caught a silent indexing error before it propagated
into the main results.

\subsection{Choice of baselines}

Every reported score is accompanied by a floor, and the correct floor depends on the
question being asked. Two are used, and the difference matters when comparing against
published work:

\begin{itemize}
  \item \textbf{Most-frequent baseline} --- always predicts the single most common label
    set. This is the appropriate floor for exact-match and Hamming loss, where it is
    surprisingly hard to beat.
  \item \textbf{Frequency-based random baseline} --- predicts each genre independently at
    random with probability equal to its base rate. Its expected macro-F1 is approximately
    the mean base rate, which makes it a much higher and more honest floor for macro-F1.
\end{itemize}

\citet{ma2021computational} report the second of these (macro-F1 $0.32$ on their corpus)
and the first separately (macro-F1 $0.19$). Comparing our results against the wrong one of
the two would misstate the margin by more than a tenth of an F1 point, so both are
computed on both corpora.

\subsection{Data splitting: grouping by film}

The corpus contains 346 clips drawn from only 43 films, with several clips per film. A
random split therefore places clips of the \emph{same} film in both the training and the
test set, and since clips from one film share instrumentation, orchestration and
recording characteristics --- and share their genre label by construction --- a model can
score well by recognising the film rather than by learning anything about genre.

All reported results consequently use \textbf{GroupKFold cross-validation grouped by
film}, so that every clip of a given film falls entirely on one side of the split. The
size of the effect this controls was measured directly rather than assumed: evaluated
with an ungrouped split, the AST baseline reaches a macro-F1 of $0.44$, against $0.26$
under grouping. Roughly forty percent of its apparent performance is film memorisation.
The emotion features lose far less ($+0.05$), which is itself a substantive finding and is
discussed in Chapter~\ref{ch:evaluation}. Ungrouped splits are reported \emph{only} in
that comparison, explicitly labelled as leaky, and never as a result.

\subsection{Repeated cross-validation and hyperparameter selection}

A single five-fold split yields five scores, which is far too few to resolve differences
of the size at stake here (a few hundredths of macro-F1); early versions of this analysis
consequently mislabelled real and consistent effects as ``within noise'', which is a
statement about the resolution of the experiment rather than about the models. The
protocol is therefore \textbf{$10\times5$ repeated GroupKFold}: the assignment of films to
folds is re-randomised in each of ten repeats, giving 50 leakage-safe folds. Every system
being compared is scored on the \emph{same} folds, so all comparisons are properly paired.

Hyperparameters are selected by \textbf{nested cross-validation}. The inverse
regularisation strength $C$ of the logistic regression is chosen by an inner GroupKFold
\emph{within each training fold}, so the test fold never participates in the choice.
The inner loop is grouped by film as well: tuned on an ungrouped inner split, $C$ is
selected against film-identity leakage and comes out systematically too weak. This step
is not cosmetic --- sklearn's default $C=1.0$ turned out to be the worst setting for every
feature set, and it handicapped the high-dimensional embedding baselines more than the
eleven-dimensional emotion features, i.e.\ in the direction that would have flattered this
thesis's hypothesis. Any dimensionality reduction (PCA) and the Stage-1 emotion regressor
are likewise fitted inside the training fold only.

\subsection{Significance testing}

Two statistics accompany every headline comparison.

Confidence intervals are computed over the \textbf{ten per-repeat means} rather than over
the 50 individual folds, because folds within a repeat share training data and are not
independent observations.

For paired comparisons, a naive $t$-test over the 50 folds is anti-conservative for the
same reason: resampled folds overlap, so the variance of the difference is underestimated
and $p$-values come out too small. The \textbf{corrected resampled $t$-test} of
\citet{nadeau2003inference} is used instead, which inflates the variance estimate by the
train/test overlap factor. A nonparametric \textbf{win rate} --- the fraction of the 50
folds on which one system beats the other --- is reported alongside it, since it makes no
distributional assumption and is easy to interpret.

One further quantity is reported for non-significant comparisons: the limiting $p$-value
$p_{\text{limit}}$ that the corrected test converges to as the number of repeats grows
without bound. Only the $1/n$ term of the corrected variance shrinks with more repeats;
the overlap term does not. When $p_{\text{limit}} > 0.05$, no amount of additional
resampling can make the comparison significant and only a larger corpus can --- which
distinguishes a genuinely underpowered corpus from an under-resampled experiment.

\subsection{Cross-corpus evaluation}

The transfer experiment (Section~\ref{sec:zero_shot}) replaces cross-validation with a
strict train-on-one, test-on-the-other design: models are fitted on the Eerola corpus
alone and applied to all 110 films of the Blockbuster corpus, whose labels are used only
for scoring. Because the two corpora annotate different genre inventories, both are first
relabelled into a single shared six-genre space, so that the classifier being transferred
predicts exactly the same classes it was trained on. Where a Stage-2 classifier is trained
on \emph{predicted} emotions, those predictions are produced out-of-fold on the training
corpus, so that the classifier is fitted on the same noisy quantity it will encounter at
test time rather than on near-perfect in-sample predictions.

\subsection{Reproducibility}

All random seeds are fixed at 42; audio embeddings are extracted once, cached per clip and
version-controlled, so every result is reproducible from a clone without re-running any
neural network over the audio. Each experiment is a self-contained script, and the
headline experiments write their complete per-fold scores to machine-readable JSON so that
reported aggregates can be re-derived independently of the prose.
